\section{Zero-residence concentration and the original pointwise source}
\label{sec:v85-zero-residence}
The complete band estimate identifies exactly what the positive
residence cutoff leaves behind. Its residual is a single monotone
charge on the actual surface, and it has a direct comparison with
an original collision-source restriction. Neither construction changes
the exact event or assumes a surface spectral theorem.

For nonnegative source families put
\[
 \mathcal H_M(f)=\limsup_m\sup_{R,n,k}m^2
       \esssup_{t\in\mathcal C_{m,R,n,k}(M)} f_{m,R,n,k}(t).
\]
At finite count define the short-segment occupancy trace
\begin{equation}\label{eq:v85-zero-residence-source}
 \mathcal Z_\delta(t)=\int_{\Gamma_{\rm ng}^+}
           \one_{\{0<l\le\delta\}}\one_{\{a<t<b\}}\,d\zeta(y),
 \qquad
 \mathfrak B_M=\lim_{\delta\downarrow0}\mathcal H_M(\mathcal Z_\delta).
\end{equation}
The limit exists monotonically in $[0,\infty]$. The complete residence
tail is inside the height functional. At fixed count,
$\mathcal Z_\delta\le2\overline{\mathcal F}_{2\delta}$, so it
is locally finite almost everywhere. Such fixed-count finiteness
alone does not determine $\mathfrak B_M$.

\begin{theorem}[The exact limiting surface charge]
\label{thm:v85-boundary-charge}
For every sufficiently small fixed $r>0$,
\begin{equation}\label{eq:v85-boundary-charge}
 \mathfrak B_M\le\mathfrak F_M^{\rm all}(r)
             \le\mathfrak B_M+C_\eta r^{1/32}.
\end{equation}
In particular,
\begin{equation}\label{eq:v85-boundary-limit}
       \lim_{r\downarrow0}\mathfrak F_M^{\rm all}(r)=\mathfrak B_M.
\end{equation}
Thus full surface-height decay is equivalent to zero charge of
vanishing-residence segments. Every positive residence band has
already been estimated in Theorem~\ref{thm:v85-band-height}.
\end{theorem}
\begin{proof}
For $\delta<r$, a selected short segment satisfies $b-t<l\le\delta$
whenever it contains $t$. Therefore
\[
 (1-\delta/r)\mathcal Z_\delta(t)
       \le\overline{\mathcal F}_r(t),\qquad
 \overline{\mathcal F}_r(t)
       \le\mathcal Z_\delta(t)+\overline{\mathcal F}_r^{\ge\delta}(t).
\]
A possible equality face $l=\delta$ on the upper side only enlarges
the comparison. Apply $\mathcal H_M$ and
Theorem~\ref{thm:v85-cutoff-surface}, first fixing $r,\delta$.
Only afterwards let $\delta\downarrow0$. This gives
\eqref{eq:v85-boundary-charge}, including infinite charges by monotone
inequalities. Its squeeze gives \eqref{eq:v85-boundary-limit}.
No statement about small mass on a shrinking interval is substituted
for the essential supremum.
\end{proof}

\subsection{A comparison on the original collision source}
Let $\Omega_+^{\rm ng}$ consist of the regular states whose maximal
forward segment has a nongrazing outgoing face. On a maximal segment
put $l(x)=\tau^-(x)+\tau^+(x)$; this is constant along the segment.
The Borel parametrization above makes the source
\begin{equation}\label{eq:v85-tube-source}
 q_\delta(t)\,dt
   =F_\#\bigl(\one_{\Omega_+^{\rm ng}}\one_{\{l\le\delta\}}
                  \bar\nu\bigr),\qquad
 \mathfrak Q_M=\lim_{\delta\downarrow0}\mathcal H_M(q_\delta)
\end{equation}
unambiguous. This is the original volume measure restricted to actual
short maximal segments, with its sole original normalization. It is
not a boundary collar. One has $0\le q_\delta\le e_\delta\le p$,
where $e_\delta$ is the true-exit starting density of
Theorem~\ref{thm:v83-exit-local-mass}. In particular its fixed-width
mass is at most $CL\delta^{1/32}$ in ordered normalization; this
is not a proof that $\mathfrak Q_M=0$.

For the actual fractional source retain $0<\eta<1/2$ and define
\begin{equation}\label{eq:v85-volume-charge}
 \mathfrak A_{\eta,M}=\lim_{h\downarrow0}
                  \mathcal H_M(\mathsf P_{\eta,h}v^\eta).
\end{equation}
This is the same canonical microscopic charge as in the preceding
revision, with the offset convention \eqref{eq:v83-dyadic-tail}.

\begin{theorem}[Surface trace versus actual short-segment source]
\label{thm:v85-trace-volume}
For every fixed $\delta>0$, almost everywhere,
\begin{equation}\label{eq:v85-trace-volume}
 |\mathcal Z_\delta(t)-q_\delta(t)|
                      \le C_\eta\mathsf P_{\eta,\delta}v^\eta(t).
\end{equation}
Consequently, as extended nonnegative bounds,
\begin{equation}\label{eq:v85-charge-volume}
 \mathfrak B_M\le\mathfrak Q_M+C_\eta\mathfrak A_{\eta,M},
 \qquad
 \mathfrak Q_M\le\mathfrak B_M+C_\eta\mathfrak A_{\eta,M}.
\end{equation}
If the fractional concentration charge is zero, the actual surface
charge equals the actual short-segment collision-source charge.
\end{theorem}
\begin{proof}
On one selected segment, the two fiber expressions use, respectively,
$A_j(b_j-,\theta)$ and $A_j(t,\theta)$ with the same indicator
$a_j<t<b_j$, $l_j\le\delta$. Absolute variation gives
\[
 |A_j(b_j-,\theta)-A_j(t,\theta)|
       \le C\int_t^{b_j} A_j(u,\theta)I_m(u,\theta)/z(u,\theta)\,du.
\]
At every intermediate target $u=t+s$, radial growth yields
$z(u)\ge\sqrt{2\lambda_\circ s}$, and hence
$z(u)^{-1}\le C_\eta s^{-(1-\eta)/2}z(u)^{-\eta}$.
Theorem~\ref{thm:v85-segment-map} identifies the selected fiber sum
at $u$ with a submeasure of the same actual source $v^\eta(u)$.
The length to its first exit is at most $\delta$. Tonelli, first on
compact regular subsegments and then by positive convergence, proves
\eqref{eq:v85-trace-volume}. No inverse Jacobian is deleted: both
physical angular densities and their full variation are retained.
Local finiteness of both terms was proved above, so the difference
does not subtract two undefined infinite assignments. Apply
$\mathcal H_M$ to the two resulting positive inequalities and only
then let $\delta\downarrow0$. This proves
\eqref{eq:v85-charge-volume}.
\end{proof}

\subsection{The ordered pointwise theorem}
For clarity, the charge identities used below have their full width
domain. With $\kappa=(1-\eta)/2$ and $0<r,2^{-j}r<r_0$,
\begin{align}
 \mathfrak T_{\eta,M}(2^{-j}r)&=\mathfrak T_{\eta,M}(r)
                   &&(j\in\mathbb Z),\notag\\
 \mathfrak T_{\eta,M}(r)&\le\mathfrak A_{\eta,M}
                    \le2^\kappa\mathfrak T_{\eta,M}(r),\notag\\
 \mathcal H_M(\mathsf P_{\eta,r}v^\eta)
       &\le C_\eta r^{1-\kappa}+2^\kappa\mathfrak A_{\eta,M}.
                                             \label{eq:v85-charge-domain}
\end{align}
Indeed the dyadic annuli partition $(0,2^{-N}r]$ at finite count.
Apply the height functional to their comparison with the truncated
potential before $N\to\infty$. The identity
$T_{\eta,r/2,N}=T_{\eta,r,N+1}$ proves the first assertion on the
displayed domain. A finite head is bounded using
\eqref{eq:v83-fixed-width-input}; only its complete remaining tail
is then passed to $N\to\infty$. These are also direct proofs of
\eqref{eq:v85-charge-domain}.

\begin{theorem}[The original endpoint after all positive residence scales]
\label{thm:v85-endpoint}
For all sufficiently large fixed $B$, fixed $0<\eta<1/2$ and small
previously fixed $r>0$,
\begin{align}
 \mathcal E_M\le{}&C_MB^{-1/192}+C\mathfrak B_M
                  +C_\eta\mathfrak A_{\eta,M}\notag\\
 &+C_\eta r^{1/32}
                  +C_\Theta(1+\varepsilon(B)^{-1})\sqrt r.
                                             \label{eq:v85-endpoint}
\end{align}
One may replace $C\mathfrak B_M$ by $C\mathfrak Q_M$ after enlarging
the coefficient of $\mathfrak A_{\eta,M}$. Zero fractional charge and
zero short-segment source charge therefore suffice for the unchanged
pointwise assertion \eqref{eq:v76-target}. Their vanishing is not
asserted by the cutoff-uniform bounds proved in this article.
\end{theorem}
\begin{proof}
Use the source-level positive-remainder comparison in the proof of
Theorem~\ref{thm:v83-endpoint} before separating its tail. Insert
\eqref{eq:v85-boundary-charge} and \eqref{eq:v85-charge-domain}.
Since $1-\kappa>1/2>1/32$, the explicit potential term is absorbed
by $C_\eta r^{1/32}$. Equation~\eqref{eq:v85-charge-volume} gives
the source form. Under the two stated zero-charge conditions choose
$r(B)=c\varepsilon(B)^\gamma$, with $\gamma>2$. For each $B$ all
widths, every positive cutoff used in the finite comparisons, and
all auxiliary bands are fixed before the collision limit. Enlarge
the auxiliary bands afterwards and take $B\to\infty$ last. The
explicit errors vanish. No zero arithmetic class is normalized.
\end{proof}

The estimates above close every fixed positive residence scale of both
short-surface mechanisms in the referee report, with constants uniform
in a subsequently chosen positive residence cutoff. They identify the
remaining full surface height with one zero-residence charge and pay
its comparison with the original volume restriction. The two remaining
nonnegative requirements in \eqref{eq:v85-endpoint} are genuine
microscopic concentration assertions. They are not proved by fixed-count
integrability, local weak norms, or exchanging a cutoff with the
collision limsup. The principal target has not been replaced by any
of those weaker statements.
